\section{Iterative Loops and the Iterator Protocol}

Loops are the control-flow structures that make repeated execution possible. A conditional statement chooses one route once. A loop creates a route that can return to an earlier point and execute a block again. This repeated route may be controlled by a condition, as in \texttt{while}, or by a traversal source, as in \texttt{for}.

Python's loop model is closely connected to its object model. A \texttt{for} loop is not merely a counter mechanism. It is a protocol-driven traversal system: Python asks an iterable object for an iterator, repeatedly asks that iterator for the next value, and stops when the iterator reports completion.

This section separates the two major loop families. The \texttt{while} loop is studied as condition-controlled indefinite iteration. The \texttt{for} loop is studied as object-driven definite traversal. From there, the section explains the difference between iterable objects and iterator objects, why iterators can become exhausted, how \texttt{break} and \texttt{continue} alter loop routes, and why Python's loop-\texttt{else} clause means normal completion rather than ordinary conditional failure.